\documentclass[11pt]{article}
\usepackage[margin=0.6in,top=0.7in,bottom=0.55in]{geometry}
\usepackage{lmodern}
\usepackage{microtype}
\usepackage{fancyhdr}
\usepackage{titlesec}
\usepackage{enumitem}
\usepackage{array}
\usepackage{tabularx}
\usepackage{booktabs}
\usepackage{lastpage}
\usepackage[hidelinks]{hyperref}

\setlength{\parindent}{0pt}
\setlength{\parskip}{0.32em}
\setlength{\headheight}{22pt}
\setlist[enumerate]{leftmargin=*,itemsep=0.22em,topsep=0.2em}
\setlist[itemize]{leftmargin=*,itemsep=0.12em,topsep=0.12em}

\titleformat{\section}{\large\bfseries\scshape}{}{0pt}{}
\titleformat{\subsection}{\normalsize\bfseries}{}{0pt}{}

\pagestyle{fancy}
\fancyhf{}
\fancyhead[L]{\footnotesize Week 5 Practice Addendum --- Teacher Guide}
\fancyhead[R]{\footnotesize Page \thepage\ of \pageref{LastPage}}
\renewcommand{\headrulewidth}{0.4pt}

\begin{document}

\begin{center}
{\LARGE\bfseries Week 5 Practice Addendum}\\[0.2em]
{\large Teacher Guide, Pacing, and Full Answer Key}\\[0.25em]
{\small Student file: \texttt{Whately\_Logic\_Week5\_Practice\_Addendum.pdf}}\\[0.1em]
{\small Place: same folder as the Week 5 student reading}\\[0.15em]
\rule{0.85\textwidth}{0.5pt}
\end{center}

\section{Why this packet exists}
Week 5 packs mood, light figure, formal validity, and several named invalid patterns into one reading. The Logic Lab (\textbf{U1L5LL}) gives thin graded transfer practice. Most students need \textbf{volume} before transfer: many short labels, many moods, many valid/invalid calls.

This addendum is \textbf{formative practice}, not a second AI-graded lab. No Reading Focus. Full key lives here so you can teach from it.

\section{What students should master (and what they should not)}
\subsection*{Must own by end of Week 5}
\begin{itemize}
\item Name A/E/I/O for ordinary sentences.
\item On a three-line syllogism: find S, P, M; write mood major--minor--conclusion.
\item Treat \textbf{Barbara, Celarent, Darii, Ferio} as the four strong \textbf{first-figure} home-base forms (names + moods + plain-English ``why'').
\item Spot invalid form, especially: undistributed middle; two negative premises; illicit process; negative premise with affirmative conclusion; particular evidence with universal conclusion.
\item Run a \textbf{form test} without yet settling all facts (Week 6).
\end{itemize}

\subsection*{Do not require}
\begin{itemize}
\item All 19 traditional valid moods, or the full mnemonic poem beyond the four first-figure names.
\item Full reduction theory (consonants in medieval names).
\item Figure numbers 2--4 as a memorized catalog (placement awareness is enough).
\item Soundness / fact test as the main skill (that is Week 6).
\end{itemize}

\section{Why Barbara, Celarent, Darii, Ferio (not ``too much'')}
These four are the classical \textbf{perfect} first-figure moods. They are enough to make ``strong form'' concrete without drowning students.

\begin{center}
\begin{tabular}{|l|c|l|}
\hline
Name & Mood & Chain idea \\
\hline
Barbara & A-A-A & All S in M; all M in P $\Rightarrow$ all S in P \\
Celarent & E-A-E & All S in M; no M in P $\Rightarrow$ no S in P \\
Darii & A-I-I & Some S in M; all M in P $\Rightarrow$ some S in P \\
Ferio & E-I-O & Some S in M; no M in P $\Rightarrow$ some S not P \\
\hline
\end{tabular}
\end{center}

\vspace{0.3em}
\noindent\textbf{First-figure layout:} major premise has M as subject (M--P); minor has M as predicate (S--M). That is the clean chain. Teaching the names is worth it because students finally have four ``this is what good looks like'' models. Spelling: \textbf{Celarent}, \textbf{Darii}, \textbf{Ferio} (not Celerent / Dario / Fario).

\section{Suggested pacing}
\begin{itemize}
\item \textbf{Day 1:} Student reading + Reading Focus discussion (not this packet).
\item \textbf{Day 2:} Addendum Parts A--C in class (pairs OK). Spot-check moods on the board.
\item \textbf{Day 3:} Part D (four strong forms) whole-class; start Part E.
\item \textbf{Day 4:} Finish E--G; homework H or I if needed.
\item \textbf{Day 5:} Part I debrief; Part J memory card; then assign graded \textbf{U1L5LL}.
\item \textbf{Later:} Lit Example U1L5LE.
\end{itemize}
If you only have two practice days: do A, B (half), D, E (half), I. Skip F--G or use as challenge.

\section{Grading options}
\begin{itemize}
\item \textbf{Recommended:} completion + accuracy spot-check (e.g.\ grade Part E odds and Part I).
\item Light points: Part B (12), Part D (name match), Part E (12), Part I (4) --- teacher choice.
\item Do \textbf{not} put this packet into the AI answer-key pipeline unless you later want a formal key file; U1L5LL remains the graded transfer product.
\end{itemize}

\section{Common confusions (teach these early)}
\begin{enumerate}
\item \textbf{Mood order:} always major premise, minor premise, conclusion --- not ``order on the page'' if the paragraph scrambles them (see C6, Part I).
\item \textbf{True conclusion $\neq$ valid form.} Preview Week 6; kill this habit now.
\item \textbf{Every invalid argument is not ``undistributed middle.''} Force the flaw bank.
\item \textbf{O vs E:} ``Some S are not P'' is O, not E.
\item \textbf{Singulars:} ``Maya is \ldots'' treated as A-style for mood (standard intro-logic classroom move).
\item \textbf{Illicit process vs undistributed middle:} middle fails to link vs a term distributed in the conclusion but not in its premise.
\item \textbf{Two negatives:} two separations from a third class do not relate the first two classes.
\end{enumerate}

\section{Relation to other Week 5 products}
\begin{itemize}
\item Reading: concepts + few models. Addendum: volume. Lab: sparse graded transfer. LE: literary application.
\item Examples here avoid the reading's whales / tyrannical orders / careful proofs / Brutus-Macbeth form demos and avoid U1L5LL's Route 9 / maples / podcasts / senators / Priya set (Part E uses roses/tulips, streaming/assemblies, federal judges instead).
\end{itemize}

\section{Sources / pedigree (for you, not for students)}
Classic first-figure illustrations modernized from standard categorical-syllogism teaching (Aristotelian first figure; medieval names Barbara, Celarent, Darii, Ferio; common invalid patterns: undistributed middle, exclusive premises, illicit major/minor). School/civic wording is original classroom adaptation.

\newpage
\section{Full answer key}

\subsection*{Part A --- Proposition warm-up}
\begin{enumerate}
\item A
\item E
\item I
\item O
\item A \quad (``Every'' = All)
\item E \quad (``None \ldots is'' = No)
\item I \quad (``At least one'' = Some)
\item O \quad (``Not every volunteer finished'' $\approx$ Some volunteers are not people who finished the training module)
\end{enumerate}

\subsection*{Part B --- Full labels + mood}
\noindent Accept reasonable synonymy in term wording (e.g.\ ``dogs'' vs ``all dogs'').

\begin{enumerate}
\item[\textbf{B1}] S: dogs \quad P: warm-blooded animals \quad M: mammals \quad Mood: \textbf{A-A-A} (Barbara shape)
\item[\textbf{B2}] S: snakes \quad P: fur-bearing animals \quad M: reptiles \quad Mood: \textbf{E-A-E} (Celarent)
\item[\textbf{B3}] S: night-shift workers \quad P: citizens with a current address on file \quad M: registered voters \quad Mood: \textbf{A-I-I} (Darii)
\item[\textbf{B4}] S: stickers on that windshield \quad P: valid dashboard displays \quad M: expired parking permits \quad Mood: \textbf{E-I-O} (Ferio)
\item[\textbf{B5}] S: fire alarms in this building \quad P: procedures that clear the stairwells \quad M: emergency drills \quad Mood: \textbf{A-A-A}
\item[\textbf{B6}] S: Jordan's lunch holds \quad P: cleared lunch account(s) \quad M: unpaid cafeteria debt(s) \quad Mood: \textbf{E-A-E}
\item[\textbf{B7}] S: pool staff members \quad P: people trained in CPR \quad M: certified lifeguards \quad Mood: \textbf{A-I-I}
\item[\textbf{B8}] S: rooms on the first floor \quad P: open student workspaces \quad M: locked server closets \quad Mood: \textbf{E-I-O}
\item[\textbf{B9}] S: Maya \quad P: student with a complete application \quad M: scholarship finalist(s) \quad Mood: \textbf{A-A-A} (singulars as A)
\item[\textbf{B10}] S: this powder \quad P: allowed competition aid \quad M: banned substance \quad Mood: \textbf{E-A-E} (conclusion negative universal/singular)
\item[\textbf{B11}] S: the tablet on Row C \quad P: (things that) must be returned before last period \quad M: tablet(s) checked out from the cart \quad Mood: \textbf{A-A-A}
\item[\textbf{B12}] S: Route 14 \quad P: (routes that) stop at the north lot \quad M: late bus route(s) \quad Mood: \textbf{E-A-E}
\end{enumerate}

\subsection*{Part C --- Mood only}
\begin{enumerate}
\item A-A-A
\item E-A-E
\item A-I-I
\item E-I-O
\item A-E-E \quad (major A, minor E, conclusion E --- valid in some figures; do not demand figure name; mood letters only here)
\item \textbf{Careful:} Conclusion is about English assignments / papers on time. Premises as printed: O then A then O. If students reorder to standard major-first: major may be treated as the premise with ``papers turned in on time'' as P. Accept \textbf{O-A-O} if they keep print order and note scramble; best teaching move: rewrite as\\
\quad Major: All of these essays are English assignments. (A)\\
\quad Minor: Some essays are not papers turned in on time. (O)\\
\quad Conclusion: Some English assignments are not papers turned in on time. (O)\\
\quad That is messy distribution-wise; use this item to teach \textbf{find conclusion first}, not as a clean validity case. For mood-as-printed: \textbf{O-A-O}.
\item A-E-E
\item O-A-O \quad (print order: O major-looking, A, O) --- actually: ``Some campus maps are not up-to-date guides'' (O), ``All campus maps are free handouts'' (A), conclusion O. Mood as printed \textbf{O-A-O}.
\item A-A-A
\item E-I-O
\end{enumerate}
\noindent\textit{Teacher tip:} If C5--C8 slow the class, skip them and return after Part D.

\subsection*{Part D --- Four strong forms}
\subsubsection*{D1}
\begin{enumerate}
\item Barbara \quad A-A-A
\item Celarent \quad E-A-E
\item Darii \quad A-I-I
\item Ferio \quad E-I-O
\end{enumerate}

\subsubsection*{D2}
\begin{enumerate}
\item Conclusion: All yearbook editors are people who verify quotations.
\item Conclusion: No bike in Rack 3 is a bike left safely overnight. \\
(or: None of the bikes in Rack 3 is left safely overnight.)
\item Minor: Some seniors are certified tutors. \\
(Any I-premise that places some seniors inside certified tutors works.)
\item Major: No peanut dessert is a safe snack for that allergy table. \\
(E premise excluding peanut desserts from the safe-snack class.)
\end{enumerate}

\subsubsection*{D3 --- sample strong answers}
\begin{enumerate}
\item First figure lines S up under M and M under P (or outside P in Celarent). If premises are granted, the conclusion only reports the link already forced --- not a new leap.
\item Darii/Ferio still \textbf{force} a particular conclusion from a universal major + particular minor. ``Some'' is not a guess; it is the largest conclusion those premises can guarantee. Overclaiming ``all'' would be the weak move.
\end{enumerate}

\subsection*{Part E --- Valid or invalid}
\begin{enumerate}
\item Mood A-A-A. \textbf{Invalid.} Undistributed middle (``plants'' never links tulips into roses).
\item Mood E-E-E. \textbf{Invalid.} Two negative premises.
\item Mood A-E-E. \textbf{Invalid.} Illicit process of the major term (``government officials'' distributed in conclusion, not properly distributed in major A premise).
\item Mood A-A-A. \textbf{Valid.} First-figure Barbara shape.
\item Mood E-A-E. \textbf{Valid.} Celarent shape.
\item Mood A-I-I but middle is ``students with community hours'' as predicate of both --- \textbf{Invalid.} Undistributed middle (both premises place groups inside community-hours students; does not force juniors into honor society). \\
\textit{Note:} If written as major ``All honor society\ldots'', minor ``Some juniors are students with community hours,'' the middle is not even shared as a linking middle in first-figure form --- still invalid.
\item Mood E-I-O. \textbf{Valid.} Ferio shape (first figure).
\item Mood A-E-E. Often taught as valid \textbf{Camestres}-like second figure if middle placement is P-side; many intro courses accept \textbf{valid}. If students say invalid from first-figure-only habits, discuss placement: major A (scored above 90 = M subject), minor E (invited = M subject) --- actually both have ``invited to review dinner'' linking differently. Walk carefully:\\
\quad Standard form attempt: All above-90 are invited. No invited is barred. Therefore no above-90 is barred. This is valid (chain of exclusion). Mark \textbf{Valid}.
\item Mood I-A-A as printed. \textbf{Invalid.} Particular premise cannot yield universal conclusion about all bus riders (over-strong conclusion). Also order is scrambled.
\item Mood E-E-E. \textbf{Invalid.} Two negative premises.
\item Mood A-A-A. \textbf{Invalid.} Undistributed middle (``ceiling fixtures'').
\item Mood A-I-I. \textbf{Valid.} Darii shape.
\end{enumerate}

\subsection*{Part F --- Same shape, different subject}
\subsubsection*{F1}
Both are A-A-A inclusion chains (first-figure style if ordered M--P / S--M). \textbf{Valid} form. Shows same shape can carry different topics.

\subsubsection*{F2}
Undistributed middle (both groups under a larger class). \textbf{Invalid.} Same bad shape, different nouns.

\subsubsection*{F3}
Two negative premises. \textbf{Invalid.} Same failure mode in biology wording and school wording.

\subsection*{Part G --- Light figure}
\begin{enumerate}
\item M subject of major, M predicate of minor. \textbf{First figure: Yes.}
\item Middle: people who stretch before practice (or ``stretch before practice'' class). M is \textbf{predicate} of both premises. \textbf{First figure: No.} (Classic undistributed-middle shape.)
\item Middle: reptiles. M subject major, M predicate minor. \textbf{First figure: Yes.} (Celarent)
\item Middle: smartphones. M subject major; M predicate minor? Minor is ``No wall clocks are smartphones'' --- M is predicate. First figure pattern for E minor is different; still M-predicate in minor. Major has M as subject. \textbf{First figure-ish layout} but mood A-E-E with this placement is \textbf{invalid} (illicit major: ``devices that use batteries'' over-extended). Mark \textbf{Invalid}; illicit process of major.
\end{enumerate}

\subsection*{Part H --- Repair (accept equivalent fixes)}
\begin{enumerate}
\item Classic undistributed middle. Repair examples: change conclusion to something not forced --- or change a premise to a real first-figure link (e.g.\ All surgeons are nurses --- false but valid Barbara if paired with All nurses work at hospital --- still may not match intent). Best student move: \textbf{explain no safe patch to ``all surgeons are nurses'' from shared workplace alone}; or weaken to ``Some people who work at the hospital are surgeons; some are nurses'' (not a syllogism repair). Full credit for clear ``shared large class does not identify subclasses.''
\item Two negatives about captains do not relate sophomores to freshmen. Explain cannot prove that conclusion; any ``repair'' must change the question.
\item Illicit major: A major does not distribute ``students with strong essays.'' Repair: weaken conclusion to something not claiming all transfer students lack strong essays --- e.g.\ conclude only that no transfer student is a \textbf{scholarship winner} (but that needs different premises). Best: show original conclusion overclaims.
\item Particular + universal cannot prove all club members barred. Repair conclusion to: \textbf{Some club members are people barred from the party} (if chain is fixed in first figure: All who skipped are barred; some club members skipped $\Rightarrow$ some club members barred).
\item Affirming the consequent / wrong direction (undistributed middle on ``events that empty the building''). Repair: cannot conclude fire drill from emptying alone; need ``All events that empty the building are fire drills'' (false) or drop the conclusion.
\end{enumerate}

\subsection*{Part I --- Buried syllogisms}
\subsubsection*{I1}
Major: Every athlete who failed the concussion screen sits out the next game. (A)\\
Minor: Rivera failed the concussion screen. (A singular)\\
Conclusion: Rivera sits out the next game. (A singular)\\
Mood: \textbf{A-A-A}. \textbf{Valid} (Barbara with singulars). Ignore snacks/popcorn distractors.

\subsubsection*{I2}
Major: Nobody without a medical pass may use the elevator before 3:00. $\approx$ No person without a medical pass is a person who may use the elevator before 3:00. (E)\\
Minor: Sam has no medical pass. (A singular about Sam)\\
Conclusion: Sam may not use the elevator before 3:00. (E/A singular negative)\\
Mood: treat as \textbf{E-A-E}. \textbf{Valid} Celarent-style. Paint smell is distractor.

\subsubsection*{I3}
Major: All students who complete six college essays are applicants ready for early decision. (A)\\
Minor: Jordan completed six college essays. (A singular)\\
Conclusion: Jordan is an applicant ready for early decision. (A singular)\\
Mood: \textbf{A-A-A}. \textbf{Valid in form} (facts/definitions may be doubtful --- Week 6 point if students raise it; form is fine).

\subsubsection*{I4}
Both premises put groups inside ``pay attention.'' Conclusion identifies the two groups. Mood A-A-A with undistributed middle. \textbf{Invalid.} Undistributed middle.

\subsection*{Part J --- Quick reference (targets)}
\begin{enumerate}
\item Mood = A/E/I/O pattern of major, minor, conclusion in that order.
\item First figure: major M--P; minor S--M; conclusion S--P.
\item Barbara A-A-A; Celarent E-A-E; Darii A-I-I; Ferio E-I-O.
\item Middle term never covers the whole class needed to link S and P.
\item Two exclusions from a third class do not relate the first two classes.
\item Form test: Does the conclusion follow from the premises?
\item Fact test: Are the premises true/justified? (Week 6)
\end{enumerate}

\newpage
\section{Board scripts (3 minutes each)}

\subsection*{Undistributed middle}
Draw one big circle (``people who stretch''). Two smaller circles inside that do not necessarily overlap (``varsity runners,'' ``managers''). Say: both inside big circle $\neq$ same small circle.

\subsection*{Two negatives}
``No freshmen are captains. No sophomores are captains.'' Both groups outside captains. They might still overlap each other or not. Nothing follows about freshmen vs sophomores.

\subsection*{Barbara chain}
Three nested circles or a chain on the board: S $\subset$ M $\subset$ P $\Rightarrow$ S $\subset$ P.

\subsection*{Celarent chain}
S $\subset$ M, and M entirely outside P $\Rightarrow$ S outside P.

\section{If students still drown}
\begin{itemize}
\item Temporary rule: only first-figure examples until Part E is stable.
\item Drop figure numbers entirely; keep only ``chain vs shared-basket.''
\item Use Part D names as the only ``valid models'' for a week, then expand.
\item Do not split the unit number; extend calendar days inside Week 5 instead.
\end{itemize}

\vspace{0.6em}
\begin{center}
\textit{End of Teacher Guide --- Week 5 Practice Addendum}
\end{center}

\end{document}
